% outlook/part_intro.tex

\chapter*{Outlook}
\phantomsection
\addcontentsline{toc}{chapter}{Outlook}

% ============================================================
\section*{What do we already know?}
% ============================================================

\begin{remark}
The previous parts developed and analyzed the AIM framework in its
present form.
\end{remark}

\begin{remark}
The foundations, geometry, states, runtime services, realization
concepts, modelling structures, optimization methods, engineering
constraints, applications, and broader interpretations have now been
established.
\end{remark}

\begin{remark}
Together, these elements form a coherent framework for
alignment-based information modelling.
\end{remark}

\begin{remark}
The present work therefore provides a complete description of AIM as it
currently exists.
\end{remark}

% ============================================================
\section*{What is still missing?}
% ============================================================

\begin{remark}
No framework is ever truly complete.
\end{remark}

\begin{remark}
New technologies emerge.

Engineering workflows evolve.

Information systems become increasingly integrated.

New requirements arise from automation, digital twins, infrastructure
management, and future transportation systems.
\end{remark}

\begin{remark}
The missing element is therefore the future perspective.
\end{remark}

\begin{remark}
The question is no longer what AIM is, but what AIM may become.
\end{remark}

% ============================================================
\section*{What comes next?}
% ============================================================

\begin{remark}
This part places AIM into a broader context and explores possible
directions for future development.
\end{remark}

\begin{remark}
The chapters that follow investigate:
\begin{itemize}
\item BIM-related perspectives,
\item alignment-centric information models,
\item future engineering workflows,
\item extensions toward PIM concepts,
\item and open research questions.
\end{itemize}
\end{remark}

\begin{remark}
The objective is not to finalize the framework, but to identify
opportunities for growth and integration.
\end{remark}

\begin{remark}
Particular attention is given to the possibility that alignment may
evolve into a first-class information entity comparable to the role
currently occupied by building-centric BIM concepts.
\end{remark}

% ============================================================
\section*{Relation to the AIM Roadmap}
% ============================================================

\begin{remark}
Within the AIM roadmap, the present part corresponds to the transition
from the present framework toward future developments.
\end{remark}

\begin{center}
\begin{tikzpicture}[
  node distance=1.4cm,
  box/.style={
    draw,
    rounded corners,
    minimum width=5.0cm,
    minimum height=0.9cm,
    align=center
  }
]

\node[box] (aim)
{Present AIM};

\node[box, below=of aim] (future)
{Future Development};

\node[box, below=of future] (bim)
{BIM, PIM, and Beyond};

\draw[->, thick] (aim) -- (future);
\draw[->, thick] (future) -- (bim);

\end{tikzpicture}
\end{center}

\begin{remark}
The progression
\[
\text{Present AIM}
\rightarrow
\text{Future Development}
\rightarrow
\text{BIM, PIM, and Beyond}
\]
marks the transition from describing the current framework toward
exploring its future evolution.
\end{remark}

\begin{remark}
The outlook therefore serves not as a conclusion, but as an invitation
to continue the development of alignment-based information modelling.
\end{remark}

% ============================================================
\begin{principle}[Open Framework Principle]
\label{princ:open-framework-principle}
Within AIM, alignment-based information modelling should be understood
as an evolving framework whose future development is expected to extend
beyond the concepts presented in this thesis.
\end{principle}

% ============================================================
\begin{summary}
Discussion explains what AIM means.

Outlook explores what AIM may become.

The present part therefore identifies future directions, broader
applications, and open questions that extend beyond the scope of the
current work while building upon its foundations.
\end{summary}
